%====================================================%
% SEZIONE 1: Studio di funzione
%====================================================%

\begin{question}[ID=1.1, subtitle={Studio di una funzione polinomiale}, difficulty=facile, tags={funzioni, studio di funzione}]
Determinare dominio, segno, intersezioni con gli assi cartesiani e parità della funzione seguente:
\begin{equation*}
    f(x)=x^2-4
\end{equation*}
\begin{solution}
\begin{enumerate}
    \item $D=\mathbb{R}$.
    \item Segno: $f(x)>0$ per $x<-2\vee x>2$; $f(x)<0$ per $-2<x<2$; $f(x)=0$ per $x=\pm2$.
    \item Intersezioni: con l'asse $x$ nei punti $(-2,0)$ e $(2,0)$; con l'asse $y$ nel punto $(0,-4)$.
    \item La funzione è pari, poiché $f(-x)=f(x)$.
\end{enumerate}
\end{solution}
\end{question}

%----------------------------------------------------%

\begin{question}[ID=1.2, subtitle={Studio di una funzione razionale}, difficulty=facile, tags={funzioni, studio di funzione}]
Determinare dominio, segno, intersezioni con gli assi cartesiani e parità della funzione seguente:
\begin{equation*}
    f(x)=\dfrac{x-2}{x+1}
\end{equation*}
\begin{solution}
\begin{enumerate}
    \item $D=\mathbb{R}\setminus\{-1\}$.
    \item Segno: $f(x)>0$ per $x<-1\vee x>2$; $f(x)<0$ per $-1<x<2$; $f(x)=0$ per $x=2$.
    \item Intersezioni: con l'asse $x$ nel punto $(2,0)$; con l'asse $y$ nel punto $(0,-2)$.
    \item La funzione non è né pari né dispari, poiché $f(-x)\neq f(x)$ e $f(-x)\neq-f(x)$.
\end{enumerate}
\end{solution}
\end{question}

%----------------------------------------------------%

\begin{question}[ID=1.3, subtitle={Studio di una funzione razionale fratta}, difficulty=media, tags={funzioni, studio di funzione}]
Determinare dominio, segno, intersezioni con gli assi cartesiani e parità della funzione seguente:
\begin{equation*}
    f(x)=\dfrac{2x-4}{x^2-9}
\end{equation*}
\begin{solution}
\begin{enumerate}
    \item $D=\mathbb{R}\setminus\{-3,3\}$.
    \item Segno: $f(x)<0$ per $x<-3\vee 2<x<3$; $f(x)>0$ per $-3<x<2\vee x>3$; $f(x)=0$ per $x=2$.
    \item Intersezioni: con l'asse $x$ nel punto $(2,0)$; con l'asse $y$ nel punto $\left(0,\frac{4}{9}\right)$.
    \item La funzione non è né pari né dispari.
\end{enumerate}
\end{solution}
\end{question}

%----------------------------------------------------%

\begin{question}[ID=1.4, subtitle={Studio di una funzione razionale pari}, difficulty=media, tags={funzioni, studio di funzione}]
Determinare dominio, segno, intersezioni con gli assi cartesiani e parità della funzione seguente:
\begin{equation*}
    f(x)=\dfrac{x^2-4}{x^2-1}
\end{equation*}
\begin{solution}
\begin{enumerate}
    \item $D=\mathbb{R}\setminus\{-1,1\}$.
    \item Segno: $f(x)>0$ per $x<-2\vee -1<x<1\vee x>2$; $f(x)<0$ per $-2<x<-1\vee 1<x<2$; $f(x)=0$ per $x=\pm2$.
    \item Intersezioni: con l'asse $x$ nei punti $(-2,0)$ e $(2,0)$; con l'asse $y$ nel punto $(0,4)$.
    \item La funzione è pari, poiché $f(-x)=f(x)$.
\end{enumerate}
\end{solution}
\end{question}

%----------------------------------------------------%

\begin{question}[ID=1.5, subtitle={Studio di una funzione razionale con più fattori}, difficulty=difficile, tags={funzioni, studio di funzione}]
Determinare dominio, segno, intersezioni con gli assi cartesiani e parità della funzione seguente:
\begin{equation*}
    f(x)=\dfrac{x^2-1}{x(x^2-4)}
\end{equation*}
\begin{solution}
\begin{enumerate}
    \item $D=\mathbb{R}\setminus\{-2,0,2\}$.
    \item Segno: $f(x)<0$ per $x<-2\vee -1<x<0\vee 1<x<2$; $f(x)>0$ per $-2<x<-1\vee 0<x<1\vee x>2$; $f(x)=0$ per $x=\pm1$.
    \item Intersezioni: con l'asse $x$ nei punti $(-1,0)$ e $(1,0)$; non esiste intersezione con l'asse $y$, poiché $x=0$ non appartiene al dominio.
    \item La funzione è dispari, poiché $f(-x)=-f(x)$.
\end{enumerate}
\end{solution}
\end{question}


%----------------------------------------------------%

\begin{question}[ID=1.6, subtitle={Studio di una funzione radicale}, difficulty=facile, tags={funzioni, studio di funzione, radicali}]
Determinare dominio, segno, intersezioni con gli assi cartesiani e parità della funzione seguente:
\begin{equation*}
    f(x)=\sqrt{x-1}
\end{equation*}
\begin{solution}
\begin{enumerate}
    \item $D=[1,+\infty)$.
    \item Segno: $f(x)\geq0$ per ogni $x\in D$; $f(x)=0$ per $x=1$.
    \item Intersezioni: con l'asse $x$ nel punto $(1,0)$; non esiste intersezione con l'asse $y$.
    \item La funzione non è né pari né dispari.
\end{enumerate}
\end{solution}
\end{question}

%----------------------------------------------------%

\begin{question}[ID=1.7, subtitle={Studio di una funzione logaritmica}, difficulty=facile, tags={funzioni, studio di funzione, logaritmi}]
Determinare dominio, segno, intersezioni con gli assi cartesiani e parità della funzione seguente:
\begin{equation*}
    f(x)=\ln x
\end{equation*}
\begin{solution}
\begin{enumerate}
    \item $D=(0,+\infty)$.
    \item Segno: $f(x)<0$ per $0<x<1$; $f(x)=0$ per $x=1$; $f(x)>0$ per $x>1$.
    \item Intersezioni: con l'asse $x$ nel punto $(1,0)$; non esiste intersezione con l'asse $y$.
    \item La funzione non è né pari né dispari.
\end{enumerate}
\end{solution}
\end{question}

%----------------------------------------------------%

\begin{question}[ID=1.8, subtitle={Studio di una funzione esponenziale}, difficulty=facile, tags={funzioni, studio di funzione, esponenziali}]
Determinare dominio, segno, intersezioni con gli assi cartesiani e parità della funzione seguente:
\begin{equation*}
    f(x)=e^x-1
\end{equation*}
\begin{solution}
\begin{enumerate}
    \item $D=\mathbb{R}$.
    \item Segno: $f(x)<0$ per $x<0$; $f(x)=0$ per $x=0$; $f(x)>0$ per $x>0$.
    \item Intersezioni: con l'asse $x$ e con l'asse $y$ nel punto $(0,0)$.
    \item La funzione non è né pari né dispari.
\end{enumerate}
\end{solution}
\end{question}

%----------------------------------------------------%

\begin{question}[ID=1.9, subtitle={Studio di una funzione esponenziale}, difficulty=facile, tags={funzioni, studio di funzione, esponenziali}]
Determinare dominio, segno, intersezioni con gli assi cartesiani e parità della funzione seguente:
\begin{equation*}
    f(x)=2^x
\end{equation*}
\begin{solution}
\begin{enumerate}
    \item $D=\mathbb{R}$.
    \item Segno: $f(x)>0$ per ogni $x\in\mathbb{R}$.
    \item Non esistono intersezioni con l'asse $x$; con l'asse $y$ nel punto $(0,1)$.
    \item La funzione non è né pari né dispari.
\end{enumerate}
\end{solution}
\end{question}

%----------------------------------------------------%

\begin{question}[ID=1.10, subtitle={Studio di una funzione radicale pari}, difficulty=media, tags={funzioni, studio di funzione, radicali, parità}]
Determinare dominio, segno, intersezioni con gli assi cartesiani e parità della funzione seguente:
\begin{equation*}
    f(x)=\sqrt{x^2-4}
\end{equation*}
\begin{solution}
\begin{enumerate}
    \item $D=(-\infty,-2]\cup[2,+\infty)$.
    \item Segno: $f(x)\geq0$ per ogni $x\in D$; $f(x)=0$ per $x=\pm2$.
    \item Intersezioni: con l'asse $x$ nei punti $(-2,0)$ e $(2,0)$; non esiste intersezione con l'asse $y$.
    \item La funzione è pari, poiché $f(-x)=f(x)$.
\end{enumerate}
\end{solution}
\end{question}

%----------------------------------------------------%

\begin{question}[ID=1.11, subtitle={Studio di una funzione logaritmica pari}, difficulty=media, tags={funzioni, studio di funzione, logaritmi, parità}]
Determinare dominio, segno, intersezioni con gli assi cartesiani e parità della funzione seguente:
\begin{equation*}
    f(x)=\ln(x^2+1)
\end{equation*}
\begin{solution}
\begin{enumerate}
    \item $D=\mathbb{R}$.
    \item Segno: $f(x)\geq0$ per ogni $x\in\mathbb{R}$; $f(x)=0$ solamente per $x=0$.
    \item Intersezioni: con l'asse $x$ e con l'asse $y$ nel punto $(0,0)$.
    \item La funzione è pari, poiché $f(-x)=f(x)$.
\end{enumerate}
\end{solution}
\end{question}

%----------------------------------------------------%

\begin{question}[ID=1.12, subtitle={Studio di una funzione logaritmica con denominatore}, difficulty=media, tags={funzioni, studio di funzione, logaritmi}]
Determinare dominio, segno, intersezioni con gli assi cartesiani e parità della funzione seguente:
\begin{equation*}
    f(x)=\frac{\ln x}{x}
\end{equation*}
\begin{solution}
\begin{enumerate}
    \item $D=(0,+\infty)$.
    \item Segno: $f(x)<0$ per $0<x<1$; $f(x)=0$ per $x=1$; $f(x)>0$ per $x>1$.
    \item Intersezioni: con l'asse $x$ nel punto $(1,0)$; non esiste intersezione con l'asse $y$.
    \item La funzione non è né pari né dispari.
\end{enumerate}
\end{solution}
\end{question}

%----------------------------------------------------%

\begin{question}[ID=1.13, subtitle={Studio di una funzione radicale fratta}, difficulty=difficile, tags={funzioni, studio di funzione, radicali}]
Determinare dominio, segno, intersezioni con gli assi cartesiani e parità della funzione seguente:
\begin{equation*}
    f(x)=\sqrt{\frac{x-1}{x+1}}
\end{equation*}
\begin{solution}
\begin{enumerate}
    \item $D=(-\infty,-1)\cup[1,+\infty)$.
    \item Segno: $f(x)\geq0$ per ogni $x\in D$; $f(x)=0$ per $x=1$.
    \item Intersezioni: con l'asse $x$ nel punto $(1,0)$; non esiste intersezione con l'asse $y$.
    \item La funzione non è né pari né dispari.
\end{enumerate}
\end{solution}
\end{question}

%----------------------------------------------------%

\begin{question}[ID=1.14, subtitle={Studio di una funzione esponenziale razionale}, difficulty=difficile, tags={funzioni, studio di funzione, esponenziali, parità}]
Determinare dominio, segno, intersezioni con gli assi cartesiani e parità della funzione seguente:
\begin{equation*}
    f(x)=\frac{e^x-1}{e^x+1}
\end{equation*}
\begin{solution}
\begin{enumerate}
    \item $D=\mathbb{R}$.
    \item Segno: $f(x)<0$ per $x<0$; $f(x)=0$ per $x=0$; $f(x)>0$ per $x>0$.
    \item Intersezioni: con l'asse $x$ e con l'asse $y$ nel punto $(0,0)$.
    \item La funzione è dispari, poiché $f(-x)=-f(x)$.
\end{enumerate}
\end{solution}
\end{question}

%----------------------------------------------------%

\begin{question}[ID=1.15, subtitle={Studio di una funzione logaritmica composta}, difficulty=difficile, tags={funzioni, studio di funzione, logaritmi, parità}]
Determinare dominio, segno, intersezioni con gli assi cartesiani e parità della funzione seguente:
\begin{equation*}
    f(x)=\ln\left(\frac{x-1}{x+1}\right)
\end{equation*}
\begin{solution}
\begin{enumerate}
    \item $D=(-\infty,-1)\cup(1,+\infty)$.
    \item Segno: $f(x)>0$ per $x<-1$; $f(x)<0$ per $x>1$; non esistono zeri.
    \item Non esistono intersezioni né con l'asse $x$ né con l'asse $y$.
    \item La funzione è dispari, poiché $f(-x)=-f(x)$.
\end{enumerate}
\end{solution}
\end{question}


%====================================================%
% SEZIONE 2: Studio completo di funzione
%====================================================%

\begin{question}[ID=2.1, subtitle={Studio completo di una funzione polinomiale}, difficulty=facile, tags={funzioni, studio di funzione, polinomi}]
Effettua lo studio completo della funzione seguente. Determina dominio, intersezioni con gli assi cartesiani, segno, parità, limiti, eventuali asintoti, derivata prima, intervalli di monotonia, eventuali estremi relativi, derivata seconda e concavità. Rappresenta infine schematicamente il grafico.
\begin{equation*}
    f(x)=x^2-4x+3
\end{equation*}
\begin{solution}
\begin{enumerate}
    \item $D=\mathbb{R}$.
    \item Intersezioni: con l'asse $x$ nei punti $(1,0)$ e $(3,0)$; con l'asse $y$ nel punto $(0,3)$.
    \item Segno: $f(x)>0$ per $x<1\vee x>3$; $f(x)<0$ per $1<x<3$; $f(x)=0$ per $x=1,3$.
    \item La funzione non è né pari né dispari; presenta simmetria rispetto alla retta $x=2$.
    \item $\displaystyle \lim_{x\to\pm\infty}f(x)=+\infty$.
    \item Non presenta asintoti.
    \item $f'(x)=2x-4=2(x-2)$.
    \item La funzione è decrescente per $x<2$ e crescente per $x>2$.
    \item In $x=2$ si ha un minimo relativo e assoluto: $f(2)=-1$.
    \item $f''(x)=2>0$: la funzione è convessa su tutto $\mathbb{R}$ e non presenta punti di flesso.
\end{enumerate}
\end{solution}
\end{question}

%----------------------------------------------------%

\begin{question}[ID=2.2, subtitle={Studio completo di una funzione razionale}, difficulty=media, tags={funzioni, studio di funzione, funzioni razionali}]
Effettua lo studio completo della funzione seguente. Determina dominio, intersezioni con gli assi cartesiani, segno, parità, limiti, eventuali asintoti, derivata prima, intervalli di monotonia, eventuali estremi relativi, derivata seconda e concavità. Rappresenta infine schematicamente il grafico.
\begin{equation*}
    f(x)=\dfrac{x-1}{x+1}
\end{equation*}
\begin{solution}
\begin{enumerate}
    \item $D=\mathbb{R}\setminus\{-1\}$.
    \item Intersezioni: con l'asse $x$ nel punto $(1,0)$; con l'asse $y$ nel punto $(0,-1)$.
    \item Segno: $f(x)>0$ per $x<-1\vee x>1$; $f(x)<0$ per $-1<x<1$; $f(x)=0$ per $x=1$.
    \item La funzione non è né pari né dispari.
    \item $\displaystyle \lim_{x\to-1^-}f(x)=+\infty$ e $\displaystyle \lim_{x\to-1^+}f(x)=-\infty$.
    \item $\displaystyle \lim_{x\to\pm\infty}f(x)=1$.
    \item Gli asintoti sono $x=-1$ verticale e $y=1$ orizzontale.
    \item $f'(x)=\dfrac{2}{(x+1)^2}>0$ per ogni $x\neq-1$.
    \item La funzione è crescente negli intervalli $(-\infty,-1)$ e $(-1,+\infty)$ e non presenta estremi relativi.
    \item $f''(x)=-\dfrac{4}{(x+1)^3}$; la funzione è convessa per $x<-1$ e concava per $x>-1$. Non presenta punti di flesso perché $x=-1$ non appartiene al dominio.
\end{enumerate}
\end{solution}
\end{question}

%----------------------------------------------------%

\begin{question}[ID=2.3, subtitle={Studio completo di una funzione razionale fratta}, difficulty=media, tags={funzioni, studio di funzione, funzioni razionali}]
Effettua lo studio completo della funzione seguente. Determina dominio, intersezioni con gli assi cartesiani, segno, parità, limiti, eventuali asintoti, derivata prima, intervalli di monotonia, eventuali estremi relativi, derivata seconda e concavità. Rappresenta infine schematicamente il grafico.
\begin{equation*}
    f(x)=\dfrac{x^2-1}{x-2}
\end{equation*}
\begin{solution}
\begin{enumerate}
    \item $D=\mathbb{R}\setminus\{2\}$.
    \item Intersezioni: con l'asse $x$ nei punti $(-1,0)$ e $(1,0)$; con l'asse $y$ nel punto $\left(0,\tfrac12\right)$.
    \item Segno: $f(x)<0$ per $x<-1\vee 1<x<2$; $f(x)>0$ per $-1<x<1\vee x>2$; $f(x)=0$ per $x=\pm1$.
    \item La funzione non è né pari né dispari.
    \item $\displaystyle \lim_{x\to2^-}f(x)=-\infty$ e $\displaystyle \lim_{x\to2^+}f(x)=+\infty$.
    \item $\displaystyle \lim_{x\to\pm\infty}f(x)=\pm\infty$ e
    \[
    f(x)=x+2+\dfrac{3}{x-2}.
    \]
    Quindi l'asintoto verticale è $x=2$ e quello obliquo è $y=x+2$.
    \item $f'(x)=1-\dfrac{3}{(x-2)^2}$.
    \item $f'(x)=0$ per $x=2\pm\sqrt3$. La funzione è crescente per $x<2-\sqrt3$ e per $x>2+\sqrt3$, decrescente per $2-\sqrt3<x<2$ e per $2<x<2+\sqrt3$.
    \item In $x=2-\sqrt3$ si ha un massimo relativo:
    \[
    f(2-\sqrt3)=4-2\sqrt3.
    \]
    In $x=2+\sqrt3$ si ha un minimo relativo:
    \[
    f(2+\sqrt3)=4+2\sqrt3.
    \]
    \item $f''(x)=\dfrac{6}{(x-2)^3}$; la funzione è concava per $x<2$ e convessa per $x>2$. Non presenta punti di flesso perché $x=2$ non appartiene al dominio.
\end{enumerate}
\end{solution}
\end{question}

%----------------------------------------------------%

\begin{question}[ID=2.4, subtitle={Studio completo di una funzione esponenziale fratta}, difficulty=difficile, tags={funzioni, studio di funzione, esponenziali}]
Effettua lo studio completo della funzione seguente. Determina dominio, intersezioni con gli assi cartesiani, segno, parità, limiti, eventuali asintoti, derivata prima, intervalli di monotonia, eventuali estremi relativi, derivata seconda e concavità. Rappresenta infine schematicamente il grafico.
\begin{equation*}
    f(x)=\dfrac{e^x}{x}
\end{equation*}
\begin{solution}
\begin{enumerate}
    \item $D=\mathbb{R}\setminus\{0\}$.
    \item Non esistono intersezioni con l'asse $x$ né con l'asse $y$.
    \item Segno: $f(x)<0$ per $x<0$; $f(x)>0$ per $x>0$.
    \item La funzione non è né pari né dispari.
    \item $\displaystyle \lim_{x\to0^-}f(x)=-\infty$ e $\displaystyle \lim_{x\to0^+}f(x)=+\infty$.
    \item $\displaystyle \lim_{x\to+\infty}f(x)=+\infty$ e $\displaystyle \lim_{x\to-\infty}f(x)=0^-$.
    \item L'asintoto verticale è $x=0$ e l'asintoto orizzontale per $x\to-\infty$ è $y=0$.
    \item $f'(x)=\dfrac{e^x(x-1)}{x^2}$.
    \item La funzione è decrescente per $x<0$ e per $0<x<1$, crescente per $x>1$.
    \item In $x=1$ si ha un minimo relativo e assoluto:
    \[
    f(1)=e.
    \]
    \item $f''(x)=\dfrac{e^x(x^2-2x+2)}{x^3}$; poiché $x^2-2x+2=(x-1)^2+1>0$, la funzione è concava per $x<0$ e convessa per $x>0$. Non presenta punti di flesso perché $x=0$ non appartiene al dominio.
\end{enumerate}
\end{solution}
\end{question}

%----------------------------------------------------%

\begin{question}[ID=2.5, subtitle={Studio completo di una funzione logaritmica composta}, difficulty=difficile, tags={funzioni, studio di funzione, logaritmi}]
Effettua lo studio completo della funzione seguente. Determina dominio, intersezioni con gli assi cartesiani, segno, parità, limiti, eventuali asintoti, derivata prima, intervalli di monotonia, eventuali estremi relativi, derivata seconda e concavità. Rappresenta infine schematicamente il grafico.
\begin{equation*}
    f(x)=\ln\left(\dfrac{x-1}{x+1}\right)
\end{equation*}
\begin{solution}
\begin{enumerate}
    \item $D=(-\infty,-1)\cup(1,+\infty)$.
    \item Non esistono intersezioni con l'asse $x$ né con l'asse $y$.
    \item Segno: $f(x)>0$ per $x<-1$; $f(x)<0$ per $x>1$; non presenta zeri.
    \item La funzione è dispari, poiché
    \[
    f(-x)=\ln\left(\dfrac{x+1}{x-1}\right)=-f(x).
    \]
    \item $\displaystyle \lim_{x\to-1^-}f(x)=+\infty$ e $\displaystyle \lim_{x\to1^+}f(x)=-\infty$.
    \item $\displaystyle \lim_{x\to-\infty}f(x)=0^+$ e $\displaystyle \lim_{x\to+\infty}f(x)=0^-$.
    \item Gli asintoti verticali sono $x=-1$ e $x=1$; l'asintoto orizzontale è $y=0$.
    \item
    \[
    f'(x)=\frac{2}{x^2-1}>0
    \]
    per ogni $x\in D$.
    \item La funzione è crescente negli intervalli $(-\infty,-1)$ e $(1,+\infty)$ e non presenta estremi relativi.
    \item
    \[
    f''(x)=-\frac{4x}{(x^2-1)^2}.
    \]
    La funzione è convessa per $x<-1$ e concava per $x>1$. Non presenta punti di flesso perché il punto $x=0$ non appartiene al dominio.
\end{enumerate}
\end{solution}
\end{question}

%====================================================%
% SEZIONE 3: Studio di funzioni parametriche
%====================================================%


\begin{question}[ID=3.1, subtitle={Studio di $f(x)=(2x+1)/(x+k)$ con parametro $k$}, tags={studio di funzione, parametrico}, difficulty=difficile]
Considera la funzione
\begin{equation*}
    f(x) = \frac{2x+1}{x+k}
\end{equation*}
dove $k$ è un parametro reale, e indica con $y$ il suo grafico.
\begin{enumerate}
    \item Dimostra che per $k=\frac{1}{2}$ la funzione attraversa solo il primo e secondo quadrante.
    \item Studia la funzione al variare del parametro $k$ e tracciane il grafico.
    \item Posto $k=1$ rappresenta la funzione $g(x) = |f(x)|$ e determina eventuali punti di non derivabilità della funzione.
\end{enumerate}
\begin{solution}
\begin{enumerate}
\item Per $k=\tfrac12$ il denominatore è $x+\tfrac12=\tfrac{2x+1}{2}$, quindi $f(x)=\dfrac{2x+1}{(2x+1)/2}=2$ per ogni $x\neq-\tfrac12$. Il grafico è la retta orizzontale $y=2$ (privata del punto $x=-\tfrac12$): essendo $y=2>0$, attraversa solo il primo e il secondo quadrante.
\item Per $k\neq\tfrac12$: dominio $x\neq-k$; asintoto verticale $x=-k$ e asintoto orizzontale $y=2$. La derivata è $f'(x)=\dfrac{2k-1}{(x+k)^2}$: nessun estremo, funzione crescente se $k>\tfrac12$, decrescente se $k<\tfrac12$. È un'iperbole equilatera di centro $(-k,2)$ (per $k=\tfrac12$ degenera nella retta $y=2$).
\item Per $k=1$, $f(x)=\dfrac{2x+1}{x+1}$ si annulla in $x=-\tfrac12$, dove $f'=\dfrac{1}{(x+1)^2}\neq0$. Quindi $g(x)=|f(x)|$ ha in $x=-\tfrac12$ un \textbf{punto angoloso} (la curva tocca l'asse $x$ con pendenza non nulla e viene ribaltata).
\end{enumerate}
\end{solution}
\end{question}

%-------------------------------%

\begin{question}[ID=3.2, subtitle={Studio di $f(x)=x^2/(x-k)$ con parametro $k$}, tags={studio di funzione, parametrico}, difficulty=difficile]
Considera la funzione
\begin{equation*}
    f(x) = \frac{x^2}{x-k}
\end{equation*}
dove $k$ è un parametro reale, e indica con $y$ il suo grafico.
\begin{enumerate}
    \item Studia la funzione al variare del parametro $k$ e tracciane il grafico.
    \item Posto $k=0$ rappresenta la funzione $g(x) = |f(x)|$ e determina eventuali punti di non derivabilità della funzione.
\end{enumerate}
\begin{solution}
\begin{enumerate}
\item Dominio $x\neq k$. Per $k\neq0$: asintoto verticale $x=k$; dalla divisione $\dfrac{x^2}{x-k}=x+k+\dfrac{k^2}{x-k}$ segue l'asintoto obliquo $y=x+k$. Derivata $f'(x)=\dfrac{x(x-2k)}{(x-k)^2}$, che si annulla in $x=0$ e $x=2k$: per $k>0$ massimo in $(0,0)$ e minimo in $(2k,4k)$ (per $k<0$ i ruoli si scambiano).
\item Per $k=0$, $f(x)=\dfrac{x^2}{x}=x$ per $x\neq0$ (retta privata dell'origine), quindi $g(x)=|x|$ con un buco in $x=0$. Il punto angoloso di $|x|$ cadrebbe proprio in $x=0$, che però è \emph{escluso} dal dominio: nel suo dominio $g$ è perciò derivabile ovunque (nessun punto di non derivabilità effettivo).
\end{enumerate}
\end{solution}
\end{question}

%-------------------------------%

\begin{question}[ID=3.3, subtitle={Potenziale del deuterio: $V(r)=a|r|e^{1-br^2}$}, tags={studio di funzione, parametrico, fisica}, difficulty=difficile]
Il potenziale di un atomo di deuterio può essere rappresentato dalla seguente funzione:
\begin{equation*}
    V(r) = a|r|\cdot e^{1-br^2}
\end{equation*}
dove $a$ e $b$ sono due parametri reali, e $r$ la distanza dal centro.
\begin{enumerate}
    \item Determina i valori di $a$ e $b$, sapendo che il massimo valore di potenziale si trova nei punti di ascissa $1$ e $-1$ ed ha valore $\sqrt{e}$.
    \item Verificato che, in corrispondenza dei valori trovati precedentemente, l'espressione analitica della funzione è $V(r) = |r| \cdot e^{1-\frac{1}{2}r^2}$, tracciane il grafico.
    \item Verifica la presenza di un punto di non derivabilità nel centro dell'atomo e classificalo.
\end{enumerate}
\begin{solution}
$a=1,\ b=\tfrac{1}{2}$. In $r=0$: cuspide (derivata sinistra $-1$, destra $+1$ a meno del fattore $e$).
\end{solution}
\end{question}

%-------------------------------%
